% Generated draft: requires mathematical and source review.
\documentclass[11pt]{article}
\usepackage[margin=25mm]{geometry}
\usepackage{fontspec}
\setmainfont{Latin Modern Roman}
\usepackage{amsmath,amssymb,mathtools,unicode-math}
\setmathfont{Latin Modern Math}
\usepackage{enumitem,needspace,xurl,hyperref}
\setlength{\parindent}{0pt}
\setlength{\parskip}{5pt}
\setlength{\emergencystretch}{3em}
\newcommand{\N}{\mathbb{N}}
\newcommand{\Z}{\mathbb{Z}}
\newcommand{\Q}{\mathbb{Q}}
\newcommand{\R}{\mathbb{R}}
\newcommand{\C}{\mathbb{C}}
\newcommand{\F}{\mathbb{F}}
\newcommand{\A}{\mathbb{A}}
\newcommand{\PP}{\mathbb{P}}
\newcommand{\E}{\mathbb{E}}
\newcommand{\eps}{\varepsilon}
\newcommand{\dd}{\,\mathrm{d}}
\begin{document}
\section*{UnsolvedMath ID: 30000936}
\textbf{Source problem code:} catalogue entry
\subsection{Definitions and mathematical statement}
Let \(n\geq1\), and let \(F\in\Z[x_1,\ldots,x_n]\) be a nonzero form, that is, a homogeneous polynomial, of degree \(d\geq1\). The first question is to identify conditions under which there exists
\[
x\in\Z^n\setminus\{0\}\quad\text{with}\quad F(x)=0.
\]
For real \(B\geq1\), define the exact counting function from the source by
\[
N_F(B)=\#\{x\in\Z^n:F(x)=0,\quad\|x\|_\infty=\max_i|x_i|\leq B\}.
\]
The second question asks for its behavior as \(B\to\infty\), for fixed \(F\), particularly for nonsingular cubic and quartic forms beyond the diagonal or sufficiently-many-variable cases treated in the source context.

In the usual projective sense, nonsingularity means that the projective hypersurface \(F=0\) over \(\overline{\Q}\) is smooth, or equivalently that at every nonzero zero of \(F\) at least one partial derivative is nonzero. The source does not give a different nonsingularity convention. The count includes the origin, all nonprimitive solutions, both signs, and all scalar multiples within the box; it is not a count of projective points or only primitive solutions. Its nonzero-solution count is \(N_F(B)-1\).

This is a program of existence and counting questions. No specific universal main term, variable threshold, local-solubility hypothesis, positivity condition, error exponent, or uniformity as \(F\) varies is supplied, so none is added as a conjectured theorem. The desired conditions and asymptotics must depend on an explicitly specified class of forms.
\subsection{Short English statement}
Find conditions guaranteeing a nonzero integral zero of a homogeneous form, and describe the number of all integral zeros in a growing sup-norm box. Nonsingular cubic and quartic forms outside the stated special regimes are the highlighted targets.
\subsection{Sources}
\begin{itemize}
\item Local UnsolvedMath registry, permanent ID \(30000936\), source code catalogue entry, ``Integral Zeros and Counting for Polynomial Forms.'' No individual proposer is supplied.
\item \emph{Analytic Number Theory} (2008), Oberwolfach Reports, DOI 10.4171/owr/2008/14, pp. 669--746, also listed as ``Oberwolfach Report 14/2008'': \url{https://doi.org/10.4171/owr/2008/14}.
\item Supplied reference, ``Browning and Heath-Brown, forms in many variables'': \url{https://doi.org/10.4171/JEMS/668}.
\item Supplied literature triage, checked 2026-08-21: status ``partially solved.'' It reports circle-method asymptotics for many nonsingular forms in sufficiently many variables and no known general asymptotic description for arbitrary higher-degree forms. Exact hypotheses and thresholds are not reproduced. This dated assessment derives from supplied OWR curation; time-sensitive status and cited claims require independent verification.
\end{itemize}
\end{document}
